% answers_short.tex —— 2008 年·普通高等学校招生考试·上海卷·理科 · 简短答案速查
% 由 tools/make_short_answers.py 自 answers.tex 的 \gksolhead/\gkanswer 提取，
% 并与 2008/2008_数学真题_上海理_答案要点.md 逐题比对后生成，勿手改（改源 MD 后重跑）。

\begin{gknotice}
\textbf{使用说明}：本节只给最终答案便于快速核对，完整推导见解析版。答案为编写者独立推导，\textbf{非官方标准答案}。
\end{gknotice}

\gkpart{填空题}
\begin{itemize}[leftmargin=2.6em]
  \item[\textbf{1.}] $(0,2)$。
  \item[\textbf{2.}] $2$。
  \item[\textbf{3.}] $1+\mathrm i$。
  \item[\textbf{4.}] $2$。
  \item[\textbf{5.}] $\sqrt7$。
  \item[\textbf{6.}] $2$。
  \item[\textbf{7.}] $\dfrac34$。
  \item[\textbf{8.}] $(-1,0)\cup(1,+\infty)$。
  \item[\textbf{9.}] $a=10$，$b=11$。
  \item[\textbf{10.}] $h_1\cot\theta_1+h_2\cot\theta_2\leqslant2a$。
  \item[\textbf{11.}] $a<-6$ 或 $a>6$。
\end{itemize}

\gkpart{选择题}
\begin{center}\begin{tabular}{ccccc}
\toprule
题号 & 12 & 13 & 14 & 15 \\
\midrule
答案 & D & C & B & D \\
\bottomrule
\end{tabular}\end{center}

\gkpart{解答题}
\begin{itemize}[leftmargin=2.6em]
  \item[\textbf{16.}] $\arcsin\dfrac{\sqrt6}6$。
  \item[\textbf{17.}] 约 $445$ 米。
  \item[\textbf{18.}] （1）$\dfrac32$；（2）$\sqrt3$。
  \item[\textbf{19.}] （1）$\log_2(1+\sqrt2)$；（2）$m\geqslant-5$。
  \item[\textbf{20.}] （1）$(8,16)$；（2）证明见解析；（3）$P$ 的轨迹是抛物线（$y^2=2mx$ 型）。
  \item[\textbf{21.}] （1）$a_n$ 以 $1,2,3$ 循环；（2）$S_{100}=198+\dfrac{a_1}2\left(11-3^{-31}\right)$；（3）证明见解析。
\end{itemize}
